\documentclass[10pt]{article}

\usepackage[utf8]{inputenc}
\usepackage[T1]{fontenc}
\usepackage{lmodern}
\usepackage{amsmath,amssymb,amsthm}
\usepackage{array,booktabs}
\usepackage[table]{xcolor}
\usepackage{enumitem}
\usepackage{geometry}
\usepackage{hyperref}
\usepackage{algpseudocode}
\usepackage{pgfplots}
\pgfplotsset{compat=1.18}
\usepackage[ruled, vlined, algoruled]{algorithm2e}
\usepackage{tikz}
\geometry{top=0.65in,bottom=0.65in,left=0.7in,right=0.7in}
\hypersetup{
  colorlinks=true,
  linkcolor=blue!60!black,
  urlcolor=cyan!50!black,
  linktoc=all
}
\setlist{nosep,leftmargin=*}
\renewcommand{\arraystretch}{1.2}
\newcolumntype{P}[1]{>{\raggedright\arraybackslash}p{#1}}

\theoremstyle{definition}
\newtheorem{definition}{Definition}[section]
\newtheorem{theorem}[definition]{Theorem}
\newtheorem{lemma}[definition]{Lemma}
\newtheorem{remark}[definition]{Remark}

\newcommand{\OPT}{\operatorname{OPT}}
\newcommand{\ALG}{\operatorname{ALG}}
\newcommand{\poly}{\operatorname{poly}}
\newcommand{\YES}{\textsc{Yes}}
\newcommand{\NO}{\textsc{No}}
\newcommand{\NP}{\mathsf{NP}}
\newcommand{\coNP}{\mathsf{coNP}}
\newcommand{\PTIME}{\mathsf{P}}
\newcommand{\NPC}{\mathsf{NP\text{-}complete}}

\definecolor{algblue}{RGB}{0,82,155}
\definecolor{optred}{RGB}{180,45,45}
\definecolor{bridgepurple}{RGB}{105,55,145}

\title{Assignment 1}
\author{Wang Xiyu (A0282500R)}
\date{}

\begin{document}

\maketitle
\tableofcontents
\newpage

\section{MAX-CFSU}
% BEGIN ORIGINAL QUESTION 1
\subsection*{Question 1: A Collision-Free Sensor Upload [10 marks]}

A network of remote environmental sensors shares a collection of wireless gateways. During a short
emergency-upload window, the system must choose a subset of the sensors to transmit. Every
gateway $j$ can hear exactly $k\geq 2$ distinct sensors; let $B_j$ denote the set of such sensors for gateway
$j$. (Note the sets $B_j$ may overlap, i.e., for any two gateways $j\neq j'$, $B_j\cap B_{j'}$ need not be empty.)

Gateway $j$ successfully decodes a packet if \textbf{exactly one} sensor in $B_j$ transmits. If none transmits,
the gateway receives no packet; if two or more transmit, their packets collide. Successfully serving
gateway $j$ has importance $p_j\geq 0$.

The objective is to choose a set of transmitting sensors that maximizes the total importance of the
gateways that successfully decode a packet. Let $\mathrm{OPT}$ denote the maximum possible importance.

Design a randomized polynomial-time algorithm whose expected value is at least $\mathrm{OPT}/e$. Your
solution should determine and justify the transmission probability used by the algorithm, prove the
approximation guarantee, and state the running time.

\subsection*{Answer}
% END ORIGINAL QUESTION 1
Let $I$ be the sensor set and $J$ the gateway set,
For each $j\in J$, let $B_j\subseteq I$ be the sensors
it can haer, where $|B_j|=k\geq 2$,and $p_j\geq 0$ be the importance of gateway $j$.

\begin{algorithm}[H]
  \KwIn{Sensor set $I$, gateway sets $(B_j)_{j\in J}$,
        importance weights $(p_j)_{j\in J}$, probability $q$}
  \KwOut{Transmitting sensor set $X\subseteq I$}
  $X\gets\emptyset$\;
  \ForEach{$s\in I$}{
    Independently set $Z_s \gets 1$ at proabability $q$\;
    \If{$Z_s=1$}{
      $X\gets X\cup\{s\}$\;
    }
  }
  \Return{$X$}\;
\end{algorithm}
\subsection{Approximation factor}
\begin{proof}
  Define indicator binary variable $\tau$ as
  \[\tau(j) = \begin{cases}
    1, \quad j \text{ successfully decode}\\
    0, \quad \text{otherwise}
  \end{cases}\]
  For each gateway, the probability of successful transmission can be expressed as 
  \[P(j) = q \cdot (1 - q)^{k - 1} \cdot k\]
  Expectation of the importance from gateway $j$ as
  \[\mathbb{E}[p_j\tau(j)] = p_j \cdot kq(1-q)^{k-1}\]
  and the total importance of all successfully transmitted gateway as 
  \[\mathbb{E}[\text{importance}(\tau)] = \mathbb{E}[\sum_j \text{importance}(\tau(j))] = 
  \sum_j p_j\cdot kq(1-q)^{k-1}\]
Define $f(q)=kq(1-q)^{k-1}$ for $q\in[0,1]$.
Since
  \begin{align*}
  f'(q) &= k[(1 - q)^{k-1} - q(k-1)(1-q)^{k - 2}]\\
    &= k(1-kq)(1-q)^{k-2}
  \end{align*}
$f'(q)>0$ for $0<q<1/k$, and $f'(q)<0$
for $1/k<q<1$, and $f(0)=f(1)=0$, thus
the maximum occurs at $q=1/k$.\\  
  substitute $q$ with $\frac{1}{k}$, 
  \[f(q)_{\max} = (1 - \frac{1}{k})^{k - 1}\]
  \[(1 - \frac{1}{k})^{k-1} = \left(\frac{1}{1 + \frac{1}{k - 1}}\right)^{k-1} \geq \frac{1}{e}\]
  We have 
  \[\mathbb{E}[\text{importance}(\tau)] \geq \frac{1}{e} \sum_j p_j \geq \frac{1}{e} OPT\]
\end{proof}
\subsection{Algorithm}
Thus the algorithm is \\
\begin{algorithm}[H]
  \KwIn{Sensor set $I$, gateway sets $(B_j)_{j\in J}$,
        importance weights $(p_j)_{j\in J}$}
  \KwOut{Transmitting sensor set $X\subseteq I$}
  $X\gets\emptyset$\;
  \ForEach{$s\in I$}{
    Independently set $Z_s \gets 1$ at proabability $\frac{1}{k}$\;
    \If{$Z_s=1$}{
      $X\gets X\cup\{s\}$\;
    }
  }
  \Return{$X$}\;
\end{algorithm}
\subsection{Running time}
Let 
\[n=|I|, \qquad g=|J|\]
The algorithm makes $n$ independent random choices and takes
$O(n)$ time, 
including reading the gateway sets gives
$O(n+gk)$ 
\section{PBS}
% BEGIN ORIGINAL QUESTION 2
\subsection*{Question 2: Planning a Biodiversity Survey [10 marks]}

A conservation agency is choosing field-survey routes. Let $U$ be the set of species that might be
observed, and let $a_u\geq 0$ be the conservation importance of species $u\in U$. Suppose there are $m$
routes. Route $i$ costs $c_i$, observes a subset $S_i\subseteq U$, and may overlap with other routes. The agency
has budget $B$, and every route satisfies $0<c_i\leq B$. A species contributes its importance at most
once, regardless of how many chosen routes observe it.

For a set $A$ of routes, define its \textbf{value} as follows:
\[
F(A)=\sum_{u\in\bigcup_{i\in A}S_i}a_u.
\]
The goal is to maximize $F(A)$ subject to $\sum_{i\in A}c_i\leq B$.

Design a deterministic polynomial-time algorithm that returns a feasible subset of routes with value
at least $\frac12\left(1-\frac1e\right)\mathrm{OPT}$, where $\mathrm{OPT}$ is the value of an optimal budget-feasible subset. Prove the
approximation guarantee and state the running time of your algorithm.

\subsection*{Answer}
% END ORIGINAL QUESTION 2
Define marginal value of adding a new route $i$ to the collection of routes $A$  
\[\Delta(i\mid A) = F(A \cup \{i\}) - F(A)\]
and the value density per cost of route $i$ when being added to $A$ as
\[\delta(i\mid A) = \frac{\Delta (i \mid A)}{c_i}\]
\subsection{Algorithm}
\begin{algorithm}[H]
  \KwIn{Species set $U = \{u_1, u_2, \dots\}$, route set $\mathcal{S}$, budget $B$}
  \KwOut{Route collection $A$}
  Stack $A \gets \emptyset$\;
  \While{$B \geq 0 \land \mathcal{S} \neq \emptyset$ }{
    $i_{\max} = \arg\max_{i \in \mathcal{S}}\frac{\Delta(i \mid A)}{c_i}$\;
    $A.\text{push}(i_{\max})$\;
    $B \gets B - c_{i_{\max}}$\;
    $\mathcal{S} \gets \mathcal{S} \setminus \{i_{\max}\}$\;
  }
  \If{$B < 0\quad\quad$ \# this guarantees feasibility}{
    $i' \gets A.\text{top}()$\;
    $A.\text{pop}()$\;
    \Return{$\arg\max_{x \in \{A, \{i'\}\}}(F(x))$}\;
  }
  \Return $A$;
\end{algorithm}
\subsection{Approximation factor proof}
Define the following:
\begin{itemize}
  \item $\mathcal{S}$: the set of routes, $|\mathcal{S}| = m$
  \item $S_i \subseteq U$: species contaned in route $i$
  \item $A \subseteq \mathcal{S}$: chosen collection of routes
  \item $T(A) = \bigcup_{i \in A} S_i$: the set of species observed in A
  \item $F(A) = \sum_{u \in \bigcup_{i \in A}S_i} a_u$: total importance of A
  \item $c(A) = \sum_{i \in A} c_i$: total cost of A
  \item $A^*$: one optimal collection
\end{itemize}
As all importance is non negative, adding new routes never decreases total importance, meaning $F$ is 
\[OPT = F(A^*) \leq F(A \cup A^*)\]
Define the number of routes in $A^* \setminus A$ that contains some $u$ as $x_u$:
\[x_u = |\{i \in A^* \setminus A, u \in S_i\}|\]
\[u \in T(A^*) \land u \notin T(A) \implies x_u \geq 1\]
\begin{align*}
\sum_{i \in A^* \setminus A} \Delta (i \mid A) &= \sum_{i \in A^* \setminus A} [F(A\cup \{i\}) - F(A)]  \\
&= \sum_{i \in A^* \setminus A} \left[ \sum_{u \in T(A) \cup S_i} a_u - \sum_{u \in T(A)} a_u \right]\\
&= \sum_{i \in A^* \setminus A} \sum_{u \in S_i \setminus T(A)} a_u\\
&= \sum_{u \in T(A^*) \setminus T(A)} x_u a_u\\
F(A\cup A^*) - F(A) &= \sum_{u \in T(A) \cup T(A^*)} a_u - \sum_{u \in T(A)} a_u\\
&= \sum_{u \in T(A^*) \setminus T(A)} a_u\\
F(A\cup A^*) - F(A) &\leq \sum_{i \in A^* \setminus A} \Delta (i \mid A)
\end{align*}
Consider at iteration $j$, $i_j$ is selected, 
\[A_j = \{i_1, i_2, \dots, i_j\}\]
as we pick in order of the highest density, 
\[\forall i \in \mathcal{S}, \delta(i \mid A_{j-1}) \leq \delta(i_j\mid A_{j-1}) \]
corollarily 
\[\forall i \in A^* \setminus A_{j-1}, \delta(i \mid A_{j-1}) \leq \delta(i_j\mid A_{j-1}) \]
Using previous result, 
\begin{align*}
  F(A^*) - F(A_{j-1}) &\leq \sum_{i \in A^* \setminus A_{j-1}} \Delta(i \mid A_{j-1})\\
  &= \sum_{i \in A^* \setminus A_{j-1}} \frac{\Delta(i \mid A_{j-1})}{c_i} \cdot c_i\\
  &\leq \sum_{i \in A^* \setminus A_{j-1}} \frac{\Delta(i_j \mid A_{j - 1})}{c_{i_j}} \cdot c_i\\
  &= \frac{\Delta(i_j \mid A_{j - 1})}{c_{i_j}} \sum_{i \in A^* \setminus A_{j-1}} c_i
\end{align*}
From the budget constraint, 
\[\sum_{i \in A^* \setminus A} c_i \leq \sum_{i \in A^*} c_i \leq B\]
we have 
\begin{align*}
F(A^*) - F(A_{j-1}) &\leq B \cdot \frac{\Delta(i_j \mid A_{j-1})}{c_{i_j}}  \\
\Delta(i_j \mid A_{j-1}) &\geq \frac{c_{i_j}}{B}[F(A^*) - F(A_{j-1})]\\
F(A_{j}) - F(A_{j-1}) &\geq \frac{c_{i_j}}{B}[F(A^*) - F(A_{j-1})]
\end{align*}
From the definition of $\Delta$
\[F(A_j) = F(A_{j-1}) + \Delta(i_j \mid A_{j-1})\]
\begin{align*}
  F(A^*) - F(A_j) &= \left[F(A^*) - F(A_{j-1})\right] - [F(A_j) - F(A_{j-1})]\\
  &\leq \left[F(A^*) - F(A_{j-1})\right] - \frac{c_{i_j}}{B}[F(A^*) - F(A_{j-1})]\\
  &= (1 - \frac{c_{i_j}}{B})[F(A^*) - F(A_{j-1})]
\end{align*}
Let the route that causes budget overflow be $i_t$, 
Apply previous result up to $t$, and $\forall c_i, c_i \leq B$
\begin{align*}
  F(A^*) - F(A_t) &\leq (1 - \frac{c_{i_t}}{B})[F(A^*) - F(A_{t-1})]\\
  &\leq \left(1 - \frac{c_{i_t}}{B}\right)\left(1 - \frac{c_{i_{t-1}}}{B}\right) [F(A^*) - F(A_{t-2})]\\
  &\dots \\
  &\leq F(A^*)\cdot \prod_{k = 1}^t \left(1 - \frac{c_{i_k}}{B}\right)\\
  &\leq F(A^*)\cdot \prod_{k = 1}^t (e^{\frac{-c_{i_k}}{B}})\\
  &= F(A^*)\cdot e^{\left( \frac{-\sum_{k=1}^{t} c_{i_k}}{B} \right)}
\end{align*}
Since with the route added in interation $t$ included, budget constraint is violated,
\begin{align*}
\sum_{k = 1}^{t-1} c_{i_k} &\leq B < \sum_{k = 1}^{t} c_{i_k}, \quad \frac{-\sum_{k=1}^{t} c_{i_k}}{B} < -1\\ 
F(A^*) - F(A_t) &\leq F(A^*)\cdot e^{\left( \frac{-\sum_{k=1}^{t} c_{i_k}}{B} \right)} \leq F(A^*) \cdot e^{-1}\\
F(A_t) &\geq \left(1 - \frac{1}{e}\right)\cdot OPT
\end{align*}
but $F(A_t)$ is not a feasible solution, as the budget is exceeded. We will discard either the entire greedy prefix 
to keep the $i_t$ as a singleton collection, or discard $i_t$ to keep $A_{t-1}$
\[
\begin{aligned}
\mathrm{ALG}
&= \max\{F(A_{t-1}),F(\{i_t\})\} \\
&\geq \frac{1}{2}\left(F(A_{t-1})+F(\{i_t\})\right)\\
&\geq \frac{1}{2} F(A_{t-1}\cup\{i_t\}) \\
&= \frac{1}{2} F(A_t) \\
&\geq \frac{1}{2}\left(1-\frac{1}{e}\right)OPT
\end{aligned}
\]
\subsubsection*{Edge case}
If no route causes budget overflow, all routes are selected, $ALG = OPT$
\subsection{Running time}
Outer loop runs at most $|\mathcal{S}|$ times, in each iteration compute $\Delta$ takes $|U|$, finding the max among available route takes $|\mathcal{S}|$, overall
$O(|\mathcal{S}|^2|U|)$
\section{ECS}
% BEGIN ORIGINAL QUESTION 3
\subsection*{Question 3: Emergency Communication after a Storm [10 marks]}

Following a coastal storm, an emergency team may activate temporary communication relays at
candidate sites. Activating site $i$ costs $c_i\geq 0$. A community $j$ can be reached only from sites in
a specified set $N(j)$. If no activated relay reaches community $j$, the team incurs a penalty $p_j\geq 0$
for arranging a slower alternative. Each community can be reached from at most $f$ candidate sites,
where $f\geq 1$ is given; that is, $|N(j)|\leq f$ for every community $j$.

The goal is to minimize
\[
\sum_{i\text{ activated}}c_i+
\sum_{j\text{ not reached by an activated site}}p_j.
\]

\begin{enumerate}[label=(\alph*)]
  \item Formulate this problem as an integer linear program and write its natural LP relaxation. [3 marks]
  \item Using an optimal solution to the LP relaxation, give a deterministic polynomial-time rounding
  algorithm. [3 marks]
  \item Prove that your algorithm is an $(f+1)$-approximation algorithm. [4 marks]
\end{enumerate}

\subsection*{Answer}
% END ORIGINAL QUESTION 3
\subsubsection*{Goal}
Minimize 
\[\sum_{i \text{ activated}} c_i + \sum_{j \text{ not reached}} p_j\]
\subsection{Integer linear programming formulation}
Define relations $S(i)$ as site $i$ is activated; $R(j)$ as community $j$ is not reached.\\
Define selection indicator variable $x_i$ as 
\[x_i = \begin{cases}
  1, \quad S(i)\\
  0, \quad \neg S(i)
\end{cases}\quad\quad y_j = \begin{cases}
  1, \quad R(j) \\
  0, \quad \neg R(j)
\end{cases}\]
For every $j$, 
\begin{align*}
(\forall i \in N(j) \neg S(i)) \rightarrow R(j)
&\equiv \bigwedge_{i \in N(j)} \neg S(i) \rightarrow R(j)\\ 
&\equiv \sum_{i \in N(j)} (1 - x_i) - y_j \leq |N(j)| - 1\\
&\equiv |N(j)| - \sum_{i\in N(j)} x_i - y_j \leq |N(j)| - 1\\
&\equiv \sum_{i \in N(j)} x_i + y_j \geq 1
\end{align*}
Integer constraint: for every $x_i, y_j$:
\[  x_i \in \{0, 1\}\quad\quad
  y_j \in \{0, 1\}
\]
\subsubsection*{Goal}
Minimize 
\[\sum_{i} c_i x_i + \sum_{j} p_j y_j\]
\subsection{LP relaxation}
Keep the same objective and covering constraints, and replace
the binary constraints by
\[
0\leq x_i\leq 1,\qquad 0\leq y_j\leq 1.
\]
Every feasible integer solution is feasible for this relaxation,
so $LP\leq OPT$.
\subsection{Determinstic rounding} 
Let the optimal solution of the relaxed LP be $x^*_i$ for $i$, $y_j^*$ for $j$, 
Define a rounding threshold $k \in (0, \frac{1}{f})$:
\[x_i = \begin{cases}
  1, \quad x^*_i \geq k\\
  0, \quad \text{ otherwise}
\end{cases}\]
This rounding is clearly polynomial, linear to the size of site set. \\
\subsection{Approximation factor proof}
Using this deterministic rounding, in our result $\forall i, S(i) \rightarrow$
\begin{align*}
x_i^* &\geq k\\
\frac{x_i^*}{k} &\geq 1  \\
\frac{1}{k}c_i x_i^* &\geq c_i x_i \quad (x_i \leq 1)\\ 
\implies \frac{1}{k} \sum_i c_i x_i^* &\geq \sum_i c_i x_i
\end{align*}
For each $j$ that is unreached, we need to have 
\begin{align*}
\forall i \in N(j) (\neg S(i) ) \implies \forall i \in N(j) (x^*_i < k)\\
\implies \sum_{i \in N(j)} x^*_i \leq k|N(j)| \leq fk
\end{align*}
From constraint, 
\begin{align*}
  y_j^* &\geq 1 - \sum_{i\in N(j)} x_i^* \\
  y^*_j &\geq 1 - fk\\
  \frac{1}{1 - fk}\cdot y^*_j&\geq 1\\
  \frac{1}{1 - fk}\cdot p_j y^*_j &\geq p_j y_j \quad (y_j \leq 1)\\
  \frac{1}{1 - fk}\sum_j p_j y^*_j &\geq \sum_j p_j y_j\\
\end{align*}
Combining 2 parts, we have 
\[LP = \sum_i c_i x_i^* + \sum_j p_j y_j^*\]
\[ALG = \sum_i c_i x_i + \sum_j p_j y_j 
\leq \frac{1}{k} \sum_i c_i x_i^* + \frac{1}{1 - fk}\sum_j p_j y_j^* 
\leq \max\left\{\frac{1}{k}, \frac{1}{1 - fk}\right\} LP\]
\subsubsection*{Find $k$ to minimize $\max\left\{\frac{1}{k}, \frac{1}{1 - fk}\right\}$}
Notice that $g_1(k) = \frac{1}{k}$ is strictly decreasing, $g_2(k) = \frac{1}{1 - fk}$ is strictly increasing over $k \in (0, \frac{1}{f})$ ($k=0$ case is trivially impossible), for $f \geq 1$, 
$\max\left\{\frac{1}{k}, \frac{1}{1 - fk}\right\} $ occurs at the intersection when 
\begin{align*}
\frac{1}{k} &= \frac{1}{1 - fk}\\
k &= \frac{1}{f + 1}
\end{align*}
Given 
\[ALG \leq \max\left\{\frac{1}{k}, \frac{1}{1 - fk}\right\} LP\]
\[ALG \leq (f + 1) \cdot LP \leq (f + 1) \cdot OPT\]
\subsection{Running time}
Let $n$ denote the total input size, suppose the relaxed LP is solved in $O(n^a)$ time,
where $a\geq 1$ is a constant, rounding all site variables and checking which communities
are reached takes
\[O\left(|I|+|J|+\sum_{j\in J}|N(j)|\right)=O(n)\]
additional time.
Therefore, the overall running time is
$O(n^a)$, determined by the LP solver
\end{document}